\section{An Overview of Idris}

\label{sec:idris}

Idris is a purely functional programming language, with \emph{first-class}
dependent types. That is, types can be computed, passed to and returned from
functions, and stored in variables, just like any other value. In this section,
we give a brief overview of the fundamental features which we 
use in this paper. A full tutorial is available
online\footnote{\url{https://idris2.readthedocs.io/en/latest/tutorial/index.html}}.
Readers who are already familiar with Idris may skip to Section~\ref{sec:idris2}
which introduces the new implementation.

\subsection{Functions and Data Types}

\label{sec:typevars}

The syntax of Idris is heavily influenced by the syntax of Haskell. Function
application is by juxtaposition and, like Haskell and ML and other related
languages, functions are defined by recursive pattern matching equations. For
example, to append two lists:

\begin{lstlisting}
append : List a -> List a -> List a
append [] ys = ys
append (x :: xs) ys = x :: append xs ys
\end{lstlisting}

The first line is a \demph{type declaration}, which is required in
Idris\footnote{Note that unlike Haskell, we use a single colon for the type
declaration.}. Names in the type declaration which begin with a lower case
letter are \emph{type level variables}, therefore \TC{append} is parameterised
by an element type.
Data types, like \TC{List}, are defined in terms of their \demph{type
constructor} and \demph{data constructors}:

\begin{lstlisting}
data List : Type -> Type where
     Nil  : List elem
     (::) : elem -> List elem -> List elem
\end{lstlisting}

The type of types is \TC{Type}.
%\footnote{There is also an infinite hierarchy
%of cumulative universes, but the details are beyond the scope of this paper}
Therefore, this declaration states that \TC{List} is parameterised by a
\TC{Type}, and can be constructed using either \TC{Nil} (an empty list) or
\TC{::} (``cons'', a list consisting of a head element and and tail).
%
As we'll see in more detail shortly, types in Idris are first-class, thus
the type of \TC{List} (\TC{Type -> Type}) is an ordinary function type.
Syntax sugar allows us to write \texttt{[]} for \DC{Nil}, and comma separated
values in brackets expand to applications of \texttt{::}, e.g. \texttt{[1, 2]}
expands to \texttt{1 :: 2 :: Nil}.

\subsection{Interactive Programs}

\label{sec:interactive}

Idris is a pure language, therefore functions have no side effects. Like
Haskell~\cite{awkward}, we write interactive programs by \emph{describing} interactions using
a parameterised type \TC{IO}. For example, we have primitives for Console I/O, including:

\begin{lstlisting}
putStrLn : String -> IO ()
getLine : IO String
\end{lstlisting}

\TC{IO t} is the type of an interactive action which produces a result of
type \TC{t}. So, \TC{getLine} is an interactive action which, when executed,
produces a \TC{String} read from the console. Idris the language \emph{evaluates}
expressions to produce a description of an interactive action as an \TC{IO t}.
It is the job of the run time system to \emph{execute} the resulting action.
Actions can be chained using the $\mathtt{>>=}$ operator:

\begin{lstlisting}
(>>=) : IO a -> (a -> IO b) -> IO b
\end{lstlisting}

For example, to read a line from the console then echo the input:

\begin{lstlisting}
echo : IO ()
echo = getLine >>= \x => 
       putStrLn x
\end{lstlisting}

For readability, again like Haskell, Idris provides \texttt{do}-notation which
translates an imperative style syntax to applications of $\mathtt{>>=}$. The
following definition is equivalent to the definition of \texttt{echo}
above.

\begin{lstlisting}
echo : IO ()
echo = do x <- getLine
          putStrLn x
\end{lstlisting}

The translation from \texttt{do}-notation to applications of $\mathtt{>>=}$
is purely syntactic. In practice therefore we can, and do, use a more general
type for $\mathtt{>>=}$ to use \texttt{do}-notation in other contexts. We will
see this later, when implementing linear resource usage protocols.

\subsection{First-Class Types}

The main distinguishing feature of Idris compared to other languages, even
some other languages with dependent types, is that types are \emph{first-class}.
This means that we can use types anywhere that we can use any other value.
For example we can pass them as arguments to functions, return them from
functions, or store them in variables. This enables us to define type synonyms,
to compute types from data, and express relationships between and properties
of data.
% 
% \subsubsection{Type Synonyms}
% 
As an initial example, we can define a \emph{type synonym} as follows:

\begin{lstlisting}
Point : Type
Point = (Int, Int)
\end{lstlisting}

Wherever the type checker sees \FN{Point} it will evaluate it,
and treat it as \TC{(Int, Int)}:

\begin{lstlisting}
moveRight : Point -> Point
moveRight (x, y) = (x + 1, y)
\end{lstlisting}

Languages often include type synonyms as a special feature (e.g.
\texttt{typedef} in C or \texttt{type} declarations in Haskell). In Idris,
no special feature is needed. 
%First-class types also allow us to define: \emph{variadic functions}, where
%an argument is used to compute the rest of a function's type; and
%\emph{dependent data types}, where types are parameterised by values.

\subsubsection{Computing Types}

\label{sec:printf}

First-class types allow us to compute types from data.
A well-known example is \texttt{printf} in C, where
a format string determines the types of arguments to be printed.
C compilers typically use extensions to check the validity of the format string;
first-class types allow us to implement a \texttt{printf}-style variadic
function, with compile time checking of the format. We begin by defining
valid formats (limited to numbers, strings, and literals here):

\begin{lstlisting}
data Format : Type where
     Num : Format -> Format
     Str : Format -> Format
     Lit : String -> Format -> Format
     End : Format
\end{lstlisting}

This describes the expected input types. We can calculate
a corresponding function type:
  
\begin{lstlisting}
PrintfType : Format -> Type
PrintfType (Num f) = (i : Int) -> PrintfType f
PrintfType (Str f) = (str : String) -> PrintfType f
PrintfType (Lit str f) = PrintfType f
PrintfType End = String
\end{lstlisting}

A function which computes a type can be used anywhere that Idris expects a
value of type \TC{Type}. So, for the type of \FN{printf}, we name the first
argument \VV{fmt}, and use it to compute the rest of the type of \FN{printf}:

\begin{lstlisting}
printf : (fmt : Format) -> PrintfType fmt
\end{lstlisting}

We can check the type of an expression, even using an as yet undefined function,
at the Idris REPL. For example, a format corresponding to \texttt{"\%d \%s"}:

\begin{lstlisting}
Main> :t printf (Num (Lit " " (Str End)))
printf (Num (Lit " " (Str End))) : Int -> String -> String
\end{lstlisting}

We will implement this via a helper function which accumulates a string:

\begin{lstlisting}
printfFmt : (fmt : Format) -> (acc : String) -> PrintfType fmt
\end{lstlisting}

Idris has support
for \emph{interactive} program development, via text editor plugins and REPL
commands. In this paper, we will use \emph{holes} extensively.
An expression of the form \texttt{?hole} stands for an as yet unimplemented
part of a program. This defines a top level function \texttt{hole}, with a type
but no definition, which we can inspect at the REPL. So, we can write a partial
definition of \texttt{printfFmt}:

\begin{lstlisting}
printfFmt : (fmt : Format) -> (acc : String) -> PrintfType fmt
printfFmt (Num f) acc = ?printfFmt_rhs_1
printfFmt (Str f) acc = ?printfFmt_rhs_2
printfFmt (Lit str f) acc = ?printfFmt_rhs_3
printfFmt End acc = ?printfFmt_rhs_4
\end{lstlisting}

Then, if we inspect the type of \texttt{printfFmt\_rhs\_1} at the REPL, we
can see the types of the variables in scope, and the expected type of
the right hand side:

\begin{lstlisting}
Example> :t printfFmt_rhs_1
   f : Format
   acc : String
------------------------------
printfFmt_rhs_1 : Int -> PrintfType f
\end{lstlisting}

So, a format specifier of \DC{Num} means we need to write a function that
expects an \TC{Int}. 
%
% Similarly, if we inspect the type of
% \texttt{printfFmt\_rhs\_2}, we'll see that a format specifier of \DC{Str} means
% we need to write a function that expects a \TC{String}:
% 
% \begin{lstlisting}
%    f : Format
%    acc : String
% ------------------------------
% printfFmt_rhs_2 : String -> PrintfType f
% \end{lstlisting}
%
For reference, the complete definition is given in Listing~\ref{printfdef}.
As a final step (omitted here) we can write a C-style \texttt{printf} by
converting a \TC{String} to a \TC{Format} specifier.

\begin{lstlisting}[caption={The complete definition of \texttt{printf}},label=printfdef,float=h]
printfFmt : (fmt : Format) -> (acc : String) -> PrintfType fmt
printfFmt (Num f) acc = \i => printfFmt f (acc ++ show i)
printfFmt (Str f) acc = \str => printfFmt f (acc ++ str)
printfFmt (Lit str f) acc = printfFmt f (acc ++ str)
printfFmt End acc = acc
  
printf : (fmt : Format) -> PrintfType fmt
printf fmt = printfFmt fmt ""
\end{lstlisting}

This example uses data which is known at compile time---the format
specifier---to calculate the rest of the type. In Section~\ref{sec:protocols},
we will see a similar idea used to calculate a type from data which is not
known until \emph{run time}.

\subsubsection{Dependent Data Types}

We define data types (such as \TC{List a} and \TC{Format} earlier) by giving
explicit types for the \emph{type constructor} and the \emph{data
constructors}.
We defined \TC{List} as having type \TC{Type -> Type}, a function type, since
\TC{List} is parameterised by an element type. We are
not limited to parameterising by types;
type constructors can be parameterised by any value. The canonical example
is a vector, \TC{Vect}, a list with its length encoded in the type:

\begin{lstlisting}
data Vect : Nat -> Type -> Type where
     Nil  : Vect Z a
     (::) : a -> Vect k a -> Vect (S k) a
\end{lstlisting}

\TC{Nat} is the type of natural numbers, declared as follows as part of the
Idris Prelude (its standard library), where \texttt{Z} stands for ``zero'' and
\texttt{S} stands for ``successor'', here using a Haskell-style \texttt{data}
declaration, which is sugar for the fully explicit type and data constructors:

\begin{lstlisting}
data Nat = Z | S Nat
\end{lstlisting}

\textbf{Note:}
%It is convenient to use this unary declaration of numbers to
%describe lengths and sizes of other structures, because then the structure of
%\texttt{Nat} allows a \texttt{S} to correspond to adding an element. 
Idris
optimises \texttt{Nat} to a binary representation at run time, although
in practice, values of type \texttt{Nat} can often be erased when the
type system can guarantee they are not needed at run time.

Since vectors have their lengths encoded in their types, functions over vectors
encode their effect on vector length in their types:

\begin{lstlisting}
append : Vect n a -> Vect m a -> Vect (n + m) a
append [] ys = ys
append (x :: xs) ys = x :: append xs ys
\end{lstlisting}

%We will see several examples of dependent data types throughout this paper.

\subsection{Implicit Arguments}

\label{sec:implicits}

As we noted in Section~\ref{sec:typevars}, lower case names in type definitions
are type level variables. So, in the following definition\ldots

\begin{lstlisting}
append : Vect n a -> Vect m a -> Vect (n + m) a
\end{lstlisting}

\ldots \VV{n}, \VV{a} and \VV{m} are type level variables. Note that we do not
say ``type variable'' since they are not necessarily of type \TC{Type}!
Type level variables are bound as \emph{implicit arguments}. Written
out in full, the type of \FN{append} is:

\begin{lstlisting}
append : {n : Nat} -> {m : Nat} -> {a : Type} ->
         Vect n a -> Vect m a -> Vect (n + m) a
\end{lstlisting}

\textbf{Note:} We will refine this when introducing multiplicities in
Section~\ref{sec:quantities}.

The implicit arguments are concrete arguments to \FN{append}, like the
\TC{Vect} arguments. Their values are solved by
unification~\cite{millerunification,gundrythesis}.  Typically implicit
arguments are only necessary at compile time, and unused at run time.
However, we can use an implicit argument in a definition, since
the names of the arguments are in scope in the right hand side:

\begin{lstlisting}
length : {n : Nat} -> Vect n a -> Nat
length xs = n
\end{lstlisting}

One challenge with first-class types is in distinguishing those parts of
programs which are required at run time, and those which can be erased.
Tradtionally, this phase distinction has been clear: types are erased at run
time, values preserved. But this correspondence is merely a coincidence,
arising from the special (non first-class) status of types! As we see with
\FN{length} and \FN{append}, sometimes an argument might be required (in
\FN{length}) and sometimes it might be erasable (in \FN{append}).
%
Idris 1 uses a constraint solving algorithm~\cite{matusicfp}, which has been
effective in practice, but has a weakness that it is not possible to tell from
a definition's type alone which arguments are required at run time.  In
Section~\ref{sec:erasure} we will see how quantitative type theory allows us to
make a precise distinction between the run time relevant parts of a program and
the compile time only parts.

\subsection{Idris 2}

\label{sec:idris2}

Idris 2 is a new version of Idris, implemented in itself, and based on
Quantitative Type Theory (QTT) as defined in recent work by Atkey~\cite{qtt}
following initial ideas by McBride~\cite{nuttin}. 
In QTT, each variable binding is associated with a \emph{quantity} (or
\emph{multiplicity}) which denotes the number of times a variable can be
used in its scope: either zero, exactly once, or unrestricted. We will
describe these further shortly. 
Several factors have motivated the new implementation:

\begin{itemize}
\item In implementing Idris in itself, we have necessarily done the engineering
required on Idris to implement a system of this scale. Furthermore, although it is
outside the scope of the present paper, we can explore the benefits of
dependently typed programming in implementing a full-scale programming language.
\item A limitation of Idris 1 is that it is not always clear which arguments to
functions and constructors are required \emph{at run time}, and which are 
erased, even despite previous work~\cite{Brady2005,matusicfp,matus-draft}. QTT allows us
to state clearly, in a type, which arguments are erased.
Erased arguments are still relevant \emph{at compile time}.
\item There has, up to now, been no full-scale implementation of a language
based on QTT which allows exploration of the possibilities of linear and
dependent types.
\item Purely pragmatically, Idris 1 has outgrown the requirements of its
initial experimental implementation, and since significant re-engineering has
been required, it was felt that it was time to re-implement it in Idris itself.
\end{itemize}

In the following sections, we will discuss the new features in Idris 2 which
arise as a result of using QTT as the core: firstly, how to express
\emph{erasure} in the type system; and secondly, how to encode resource usage
protocols using linearity.
